\documentclass[tikz,border=0pt]{standalone}
\usepackage[fontset=none]{ctex}
\usepackage{amsmath,amssymb,bm}
\usetikzlibrary{arrows.meta,calc,decorations.pathmorphing,positioning}

\definecolor{ink}{HTML}{263238}
\definecolor{muted}{HTML}{6B7780}
\definecolor{blue}{HTML}{245BFF}
\definecolor{green}{HTML}{087F5B}
\definecolor{orange}{HTML}{E67700}
\definecolor{red}{HTML}{E03131}
\definecolor{purple}{HTML}{7048E8}
\definecolor{paleBlue}{HTML}{EDF3FF}
\definecolor{paleGreen}{HTML}{EAF8F3}
\definecolor{paleOrange}{HTML}{FFF4E6}
\definecolor{palePurple}{HTML}{F3F0FF}

\tikzset{
  title/.style={font=\fontsize{18}{22}\selectfont\bfseries,text=ink,align=center},
  subtitle/.style={font=\fontsize{13}{17}\selectfont\bfseries,text=ink,align=center},
  body/.style={font=\fontsize{10.5}{14}\selectfont,text=ink,align=center},
  small/.style={font=\fontsize{8.8}{11.5}\selectfont,text=ink,align=center},
  box/.style={draw=muted,line width=.8pt,rounded corners=2pt,inner sep=6pt,body},
  vec/.style={-{Latex[length=3.2mm,width=2.3mm]},line width=1.35pt},
  patharrow/.style={-{Latex[length=2.8mm,width=2mm]},line width=1.25pt},
  nodept/.style={circle,fill=ink,minimum size=5pt,inner sep=0pt}
}

\usepackage{fontspec,xeCJK}
\IfFontExistsTF{Heiti SC}{\setCJKmainfont[AutoFakeBold=1.5]{Heiti SC}}{\setCJKmainfont[AutoFakeBold=1.5]{FandolHei-Regular}}
\IfFontExistsTF{Heiti SC}{\setCJKsansfont{Heiti SC}}{\setCJKsansfont{FandolHei-Regular}}
\usetikzlibrary{bending}
\begin{document}
\begin{tikzpicture}[x=1cm,y=1cm]
  \path[use as bounding box] (0,0) rectangle (16,8); \fill[white] (0,0) rectangle (16,8);
  \node[title] at (8,7.45) {霍尔电压判号：先画载流子速度，再代入电荷符号};
  \draw[muted,line width=.8pt] (8,1.0)--(8,6.75);
  \node[subtitle,red] at (4,6.35) {正载流子：$q>0$};
  \node[subtitle,purple] at (12,6.35) {电子：$q<0$};
  \foreach \x in {1.1,1.9,...,7.1}{\foreach \y in {2.0,2.8,...,5.2}{
    \draw[blue] (\x,\y) circle (.08); \fill[blue] (\x,\y) circle (.028);
  }}
  \foreach \x in {9.1,9.9,...,15.1}{\foreach \y in {2.0,2.8,...,5.2}{
    \draw[blue] (\x,\y) circle (.08); \fill[blue] (\x,\y) circle (.028);
  }}
  \draw[ink,line width=1.1pt] (1.0,2.2) rectangle (7.2,5.55);
  \draw[ink,line width=1.1pt] (9.0,2.2) rectangle (15.2,5.55);
  \draw[green,vec] (1.3,3.6)--(3.0,3.6); \node[small,green,above] at (2.1,3.65) {$\mathbf j,\ \mathbf v_d$};
  \draw[green,vec] (9.3,3.6)--(11.0,3.6); \node[small,green,above] at (10.1,3.65) {$\mathbf j$};
  \draw[purple,vec] (14.9,3.0)--(13.2,3.0); \node[small,purple,below] at (14.1,2.95) {$\mathbf v_d$};
  \foreach \x in {1.5,2.2,...,6.9}{\node[small,red] at (\x,2.45) {$+$}; \node[small,blue] at (\x,5.3) {$-$};}
  \foreach \x in {9.5,10.2,...,14.9}{\node[small,blue] at (\x,2.45) {$-$}; \node[small,red] at (\x,5.3) {$+$};}
  \draw[orange,vec] (6.4,3.3)--(6.4,4.8); \node[small,orange,left] at (6.3,4.1) {$\mathbf E_H$};
  \draw[orange,vec] (14.4,4.8)--(14.4,3.3); \node[small,orange,left] at (14.3,4.1) {$\mathbf E_H$};
  \node[box,fill=paleBlue,text width=5.6cm] at (4,1.0) {$\mathbf E_H=\dfrac1{nq}\,\mathbf B\times\mathbf j$，故 $R_H=1/(nq)>0$};
  \node[box,fill=palePurple,text width=5.6cm] at (12,1.0) {$q=-e$ 时 $\mathbf v_d$ 与 $\mathbf j$ 反向，且 $R_H<0$};
\end{tikzpicture}
\end{document}
